\chapter[Sermon 20. The Lord Blames Sore the Church of Laodicea.]{The Lord Blames Sore the Church of Laodicea.}
\label{sermon:020}
\markright{Sermon 20. The Lord Blames Sore the Church of Laodicea.}

And unto the angel of the congregation which is in Laodicea, write: This saith Amen, the faithful and true witness, the beginning of the creation of God. I know your works, that you are neither cold nor hot; I would that you were either cold or hot. So then, because you are lukewarm, and neither cold nor hot, I will spit you out of my mouth. Because you say, ‘I am rich and increased with goods, and have need of nothing,’ and do not know that you are wretched, miserable, poor, blind, and naked.

\index{Repentance}The seventh and final Epistle from our Savior Christ is written by John to the Bishop of Laodicea. It is a great rebuke of that people, finding nothing worthy of praise in them. Nevertheless, it is a faithful admonition or exhortation to repentance. And following his customary practice, he indicates to whom he is writing and from whom the Epistle originates. The Epistle was commissioned by Christ to the Bishop of Laodicea and to the entire congregation. Therefore, something must be said about the Laodiceans, so that the rest may better understand and consider them.

Laodicea, the chief city of Caria, as Strabo and Pliny state, stands by the river Lycus. Antiochus Theos built the city and named it after his wife. It was the wealthiest city in Asia, as Vadiane also notes in his Epitome. It had a very profitable trade in woolen cloth. It appears that Paul also preached the gospel there. For he mentions Laodicea, and some believe he wrote his first letter to Timothy from there. Certainly, it appears that the people of Laodicea had received the gospel, but in a corrupted form. For they sought to blend the world and the church together, and to join Christ and Mammon, as is said today.

Therefore they did not abandon their greed, and their immoderate commerce (for to conduct trade moderately, without deception, no religion forbids), nor their excessive riot and pride. They did not seem to lack anything, but to possess and appear to possess all things, because they were wealthy. Against these, the Lord contends severely, declaring them to be very wretched, and more destitute than simple beggars. For just as in the church of Philadelphia he found no fault, so in this one he commends nothing at all.

You will find many today who speak like this, and it is ever on their lips: “I have learned to be a follower of the Gospel, and to be a soldier, to drink, to play the whoremonger, and to live at pleasure.” You will find churches like these, serving both Christ and Mammon, or commerce, Bacchus, Venus, and the god of battle. Both they and all such people here are refuted and called to repentance. This demonstrates that God’s mercy is greatest, not abandoning or rejecting churches so corrupt, and people so full of such great wickedness.

Woe to those who despise this boundless mercy and goodness of God and his long-suffering, and persist in their wrongdoing.

Christ is described here again most abundantly, as he is in the previous titles. Indeed, it can be gathered from all that has been said that this is the best and most complete description of Christ, so that no further explanation from human sources is needed. He reveals himself with a new name, calling him “Ho Amen,” that is, Amen. This is a Hebrew word, frequently used in the Gospels, especially in John, and by Paul in 2 Corinthians 1:20. Paul says that Christ, the Son of God, whom he and Sylvanus and Timothy preach among you, is not fickle, but in him is truth. For all the promises of God are in him, and in him they are confirmed to the praise of God through us. And the Lord explains why he calls himself that Amen: “For I am the witness,” that is, the trustworthy, faithful, constant, and true one. For Christ has been given to us by the Father to testify about the will of the Father. And his testimony, as he himself repeats often in the Gospel of John, is firm, constant, sure, certain, and true, without falsehood, doubt, or inconsistency. These things agree well with the argument in which he rebukes the Laodiceans for their sin and urges them to repentance. It is a difficult matter for the flesh to hear such doctrine; but when the certainty, assuredness, or truthfulness of the teacher is perceived, it will commonly move people’s minds, if they are not entirely hopeless and despairing.

\index{Hebrews, Epistle to the}He adds moreover another thing, which declares his dignity. For he calls himself the beginning of the creatures of God. Neither ought the Arians to seek here any defense for themselves. For it is not fitting by any one place, much less by a little word, to subvert the whole scripture, and to strive with the articles of the creed, the lively tradition of the Apostles. Our Savior Christ is considered after his deity and after his humanity. After his deity, he has no beginning, but is rather the beginning (actively as it is commonly said, not passively of all things and creatures. Neither is he a creature: for all things are made by him. Which thing both the Evangelical and Apostolic scriptures prove, John 1 and Colossians 1, and Hebrews 1, where you have passages expounding this same one. After his humanity he is called the beginning of the creature of God (namely man, which is called a creature by reason of his excellency, and for that he is the lord of creatures, for whom all things were made) as he is called the firstborn of the dead. For in Christ humankind is repaired, that it has not perished: God looked upon the countenance of his Christ when he first made man. For Christ is the beginning, that is to say, the preserver of the human nature: as it has elsewhere been told you at large. Hitherto we have had the description of Christ, which is called Amen, and the beginning of the creature of God, by whom verily all things are made, which is very and true God, witness of the divine will of God. \&c.

Now he tells the church what opinion he has of her, and what she is, that is to say, blames her. And as he has impressed upon all the former that he knows all their works, so he does to this also. First, he shows that he knows this of the church of Laodicea, and especially of its bishop, that he is neither cold nor hot. He adds, I would it were better if you were altogether cold, or thoroughly hot. But now you are lukewarm, or blood warm. An allegory taken of men’s meat, or of cold, hot, or warm water, and it is applied in a manner proverbially. He is cold, who openly follows the world, being wrapped in heathenish errors and sins of this world, and boasts nothing, or will seem to have nothing to do with the true religion. He is hot, whose breast inflames with the Holy Spirit, contemns the world, loves the true religion exceedingly, and lives a holy life. He is warm, or between both, which has neither forsaken the world, his errors, and sins, nor has fully received Christ, his truth and righteousness, but serves partly the world, partly Christ: In outward things he shows himself to be a Christian, in resorting to holy assemblies and receiving the sacraments, but inwardly he is so besieged by the world that he lives a worldly life rather than a Christian. Such a mixture the Lord does not allow, which elsewhere forbids plowing with an ox and an ass, and making a garment of linen and wool; to pour new wine into old bottles, and to patch an old garment with new cloth.

\index{Mahomet}In religious practices, a blending of different traditions is often less acceptable to God. For there are those who attempt to combine various religions, compiling many into one. Muhammad constructed his religion from Jewish and Christian teachings. Many today create a mishmash of Roman Catholicism and the Gospel, or bake a “chuchurnullis,” as the Germans call it, a cake made of diverse ingredients. If a Roman Catholic observes such a service, he does not recognize it as his own; and if an adherent of the Gospel sees it, he recognizes it as belonging to no one. For it is a mixture of the pure and the corrupt, where the pure element has no greater strength, and the corrupt element generally prevails. The masses used by many today are of this sort, neither entirely Roman Catholic nor wholly Evangelical. The Lord’s Supper does not appear in them; the Roman Catholic mass is also cut off and altered within them. If we believe that Christ established the finest rule for religion and living, why do we not follow only that Master? But we value the favor of men more, which we are unwilling to lose. We do not value the favor of Christ enough to recall the Apostle’s saying: “If I pleased men, I should not be the servant of Christ.”

But hear what the Lord says to these hypocrites. It would be better, he says, if you were either cold or hot. It would be better if you were a sinner or a heathen than a hypocrite and a mongrel. For then you might be more easily helped, according to that saying of the Lord, “If you were blind, you would have no sin.” Now, when you seem to yourselves to be just and sufficiently taught and furnished with godly doctrines and practices that please God, you leave no place for further instruction, but, despising the word of God and Christ’s institution, you prefer your mixtures before all the justifications of God. The Lord also in the Gospel says to the Pharisees: “Truly I say to you, that tax collectors and prostitutes go before you into the kingdom of God.”

The other part is clear enough; it would be better if they were fervent, namely with the Spirit of God, which thing the Apostle requests in the twelfth chapter to the Romans.

Furthermore, he threatens to plague them if they continue, as they have begun to be neutral; [illegible]. “I will spew thee out of my mouth.” By this manner of speaking, two things are signified: both the loathsomeness which God conceives of this neutrality or lukewarmness, and the vomiting out, which punishes the same. Warm water provokes a vomit. To this, he appears to have alluded, as likewise to that old phrase of speaking: “The land has vomited the Canaanites, and the same shall vomit you up also.” Therefore, these composers or mongers, with their temperature and mixture, so displease God that they engender in him a loathsomeness, are an abomination to him, so that finally he shakes them off—the same we understand of those who join together Christ and Mammon. And the phrase of speech is to be noted, “now therefore,” or “so forasmuch as,” or “now seeing it is so,” etc. Moreover, the longsuffering of God is here noted, which does not plague immediately unless there appears nowhere any hope of amendment.

He explains more fully the sin of the Laodiceans, and what causes their lukewarmness: because they love riches, in which they trust, believing they lack nothing. They consider themselves wise, and that they see all things, and are sufficiently supplied with spiritual and temporal possessions. It is less, they say, “we are rich.” More follows, “I am increased with goods”: that is to say, “I have acquired so many riches that I want nothing.”

That same he now refutes, and shows that they are utterly deceived, and are to be miserable people. For he rebukes them severely, and says, “You do not know what you are.” That ignorance is a great evil, and the beginning of desperate blindness, when a man thinks he has something that he does not. For such persevere in their error, and admit no counselor. Therefore the Lord says, “You do not know that you are poor, wretched, weary, and worn with evils.” For they are toiled with many labors that serve this world. They are miserable. You do not see your own misery; others who see it are deeply sorry. You do not see in what condition you are. This kind of speech signifies a man very wretched and desperate, whose misery others see, but he himself sees nothing, like a poor or a beggar. You think yourself very rich, but you are a stark beggar. Covetous rich men are poor; they are poor also in virtues. The people of Laodicea were blind, as the Pharisees were called blind in John 9. Well-sighted in worldly matters, yet heavily blind as betels. Naked, or destitute of good works. Void of your wedding garment. They nevertheless were richly arrayed with garments of the finest wool. But before God they appeared naked. Let the gallants of this world, or proud peacocks rather, who are so well-sighted and gorgeously appareled, mark these things well. The Lord give them understanding.
